\section{Conclusion}\label{sec:conclusion}

The viability literature under incomplete observation has been asymmetric: the sufficiency side is carried by Veliov's condition, the estimation-tube reduction and barrier certificates, while the obstruction side---why a given observation structure cannot sustain a given constraint set---has had no comparable apparatus. This paper supplies that apparatus as a calculus of certificates and shows that it is organised by a single principle: an observation-based regulator must select one response for an entire information state, and each certificate proves that the relevant selection set is empty. In finite systems the principle is exact, through a backward recursion whose one-step case is the common-action test and whose obstruction tree is a finite, mechanically checked witness. In continuous time it yields a ladder of common-action tests with a polyhedral Farkas dual, a timing bound with an explicit slack, an LP relaxation whose gap is identified with the difference between the declared and the per-step policy class, a drift certificate of finite-time exit, and a recourse certificate for the failure that appears only after information arrives. The certification layer adds the fibre criterion and the certainly-safe set as the exact limits of observation-only verdicts, and the certainty-equivalence trap as the one failure that occurs under injective observation.

Three features of the calculus are intended to be used rather than admired. Each certificate names the policy class and the adverse-path hypothesis under which it holds, so that a verdict can be traced to what it assumes. Each certificate names the response correspondence it empties, so that the failure it certifies carries its own remedy. And the gap between the certificates and the truth is located: the certificates are complete in finite systems and in the static class, jointly complete on the audited hold class, and provably incomplete as any window-measurable pair in the two-phase model. The remaining problem---a finite, checkable characterisation of information-state viability under dynamic observation---is stated with the post-observation recourse verdict as its definite target.

For monitoring and indicator design in sustainability governance the message is concrete. An observation structure must separate safety classes, not merely states; it must deliver its first informative reading within a time computable from the drift at the floor; an aggregate index can certify a floor only on its certainly-safe set and must otherwise fall silent; a known bias must be corrected inside the control law, not merely bounded; and when information arrives too late to matter, the design variable is authority after the observation rather than resolution before it. These are theorems about what no design can do, and that is the part of the design problem that has been missing.
